\section{Learning and Representation Guarantees}
\label{sec:theory}

The analysis follows the actual learning pipeline.  We first control the three
errors that arise even for a fixed frozen representation: finite-window
initialization, finite-shot measurement, and dependent finite data.  We then
resolve the representational role of width $n$ by isolating the information
lost by the shared Gaussian projection.  This width-resolved approximation
result gives positive probability to neighborhoods of good frozen designs,
which in turn makes the role of the pool size $S$ explicit.  The final theorem
combines these ingredients for the full-sample ERM predictor returned by
\cref{sec:model-selection}.

\subsection{Memory fidelity}
\label{sec:memory-fidelity}

\begin{lemma}[Exact trace-norm contraction]
\label{lem:exact-contraction}
Every branch with rate $\lambda$ satisfies
\begin{equation}
    \traceNorm{\reservoirChannel_x(\rho)-\reservoirChannel_x(\rho')}
    =\lambda\traceNorm{\rho-\rho'}\leq2\lambda.
    \label{eq:exact-trace-contraction}
\end{equation}
Hence every input history determines a unique causal state, and two
trajectories driven by the same final $w$ inputs differ by at most
$2\lambda^w$ in trace norm.
\end{lemma}

\begin{proposition}[Finite-window risk error]
\label{prop:window-certificate}
For every frozen design $\design\sim\designLaw_{\arch}$ and every
$h\in\RKHS_{m_{\arch}}^\Lambda$,
\begin{equation}
\begin{aligned}
    \abs{\risk_{\design,\infty}(h)-\risk_{\design,w}(h)}
    &\leq
    4\Lambda(\Lambda+\Upsilon)\kLip
    \left(\frac1R\sum_{r=1}^{R}\lambda_r^{2w}\right)^{1/2}\\
    &\leq
    \underbrace{4\Lambda(\Lambda+\Upsilon)\kLip
      \lambda_+^w}_{\epsWin(\arch)}.
\end{aligned}
    \label{eq:window-risk-bound}
\end{equation}
\end{proposition}

\begin{corollary}[Memory-endpoint calibration]
\label{cor:lambda-plus-calibration}
A requested tolerance $\epsilon_W>0$ is guaranteed whenever
\begin{equation}
    \lambda_0<\lambda_+\leq
    \min\!\left\{1,
    \left[\frac{\epsilon_W}
      {4\Lambda(\Lambda+\Upsilon)\kLip}
    \right]^{1/w}\right\}.
    \label{eq:lambda-plus-calibration-condition}
\end{equation}
If the right-hand upper bound is at most $\lambda_0$, the requested fidelity
is incompatible with the fixed one-step lower memory endpoint.
\end{corollary}

\paragraph{Interpretation.}
Larger $\lambda_+$ retains longer memory but slows convergence from the
finite-window implementation to its causal limit.  The calibration therefore
turns a desired memory-fidelity tolerance into an explicit admissible endpoint.
The RMS Mat\'ern geometry averages coordinatewise feature displacement, so
merely concatenating more bounded branches or observables does not worsen this
risk-sensitivity bound.  It does not claim that longer memory is always more
predictive; that question belongs to the representation term introduced below.

\subsection{Measurement fidelity}
\label{sec:measurement-fidelity}

\begin{proposition}[Grouped-measurement risk error]
\label{prop:shot-certificate}
For every frozen design and every $h\in\RKHS_{m_{\arch}}^\Lambda$,
\begin{equation}
\begin{aligned}
    \abs{\oprisk_{\design,w}(h)-\risk_{\design,w}(h)}
    &\leq
    2\Lambda(\Lambda+\Upsilon)\kLip
    \left[
      \frac1{Rq_n}\E\!\sum_{r=1}^{R}\sum_{O\in\gO_n}
      \frac{1-\phi_{r,O}(X_{-w+1:0})^2}{a_n(O)M}
    \right]^{1/2}\\
    &\leq
    2\Lambda(\Lambda+\Upsilon)\kLip
      \sqrt{\frac{1}{q_nM}\sum_{O\in\gO_n}\frac1{a_n(O)}}
      =:\epsShot(\arch,M).
\end{aligned}
    \label{eq:operational-exact-risk-bound}
\end{equation}
For the cycle cover,
$\epsShot(\arch,M)=2\Lambda(\Lambda+\Upsilon)\kLip\sqrt{5/(6M)}$.
\end{proposition}

\begin{corollary}[Measurement-budget calibration]
\label{cor:measurement-budget-calibration}
For the cycle cover, $\epsShot(\arch,M)\leq\epsilon_M$ follows from
\begin{equation}
    M\geq
    \left\lceil
      \frac{10\Lambda^2(\Lambda+\Upsilon)^2\kLip^2}{3\epsilon_M^2}
    \right\rceil.
    \label{eq:risk-measurement-budget}
\end{equation}
\end{corollary}

\paragraph{Interpretation.}
The sufficient repetitions per Pauli setting scale as
$M=O(\epsilon_M^{-2})$ with no artificial $Rn$ factor: RMS kernel geometry
prevents a larger feature vector from amplifying coordinatewise estimation
noise solely through dimension.  This does \emph{not} make a larger quantum
representation free.  The complete bank still requires $9RM$ independent
branch-state preparations/circuit runs per input window under the reusable
$n$-qubit convention, and each run uses an $n$-qubit trajectory of depth
$O(w)$ as specified in \cref{sec:topology-observables}.

\subsection{Dependent-data generalization and ERM capacity}
\label{sec:dependent-generalization}
\label{sec:dependent-penalty}

Condition on the shared preprocessor $\projection_n$ and the $S$ frozen
designs.  Let $\measurementRecord_j$ collect all candidate measurement records
for window $j$ and set
$\augprocessW_j:=(\processW_j,\measurementRecord_j)$.  Independent measurement
kernels preserve the one-step absolute-regularity coefficient, so
\begin{equation}
    \beta_{\augprocessW}(1)=\beta_{\processW}(1)
    \leq\beta_{\processZ}(g).
    \label{eq:augmented-window-mixing}
\end{equation}
For confidence $\delta\in(0,1)$, write
$\delta_g:=\delta-4(\mu-1)\beta_{\processZ}(g)$ and assume $\delta_g>0$.
Define
\begin{equation}
    \epsGen(\Lambda,S,\delta)
    :=\frac{4\Lambda(\Lambda+\Upsilon)}{\sqrt\mu}
      +3(\Lambda+\Upsilon)^2
       \sqrt{\frac{\log(4S/\delta_g)}{2\mu}}.
    \label{eq:dependent-certificate}
\end{equation}

\begin{proposition}[Uniform measured-to-operational deviation]
\label{prop:gapped-deviation}
Conditional on any realized $\projection_n$ and frozen pool, with probability
at least $1-\delta$ over the gapped sample and all training measurement
outcomes, simultaneously for every $s\leq S$ and every
$h\in\RKHS_{m_{\arch}}^\Lambda$,
\begin{equation}
    \abs{\erisk_{\design_s,w}(h)-\oprisk_{\design_s,w}(h)}
    \leq\epsGen(\Lambda,S,\delta).
    \label{eq:gapped-uniform-deviation}
\end{equation}
\end{proposition}

\begin{corollary}[Geometric-mixing gap rule]
\label{cor:geometric-mixing-gap}
Assume that the observed process has a known geometric absolute-regularity
envelope
\begin{equation}
    \beta_{\processZ}(q)\leq \mixConst e^{-\mixDecay q},
    \qquad q\geq0,
    \label{eq:geometric-mixing-envelope}
\end{equation}
for constants $\mixConst\geq1$ and $\mixDecay>0$.  For $\mu\geq2$, the choice
\begin{equation}
    g\geq
    \left\lceil
      \frac1{\mixDecay}
      \log\!\left(\frac{8\mixConst(\mu-1)}{\delta}\right)
    \right\rceil
    \label{eq:geometric-gap-rule}
\end{equation}
ensures $\delta_g\geq\delta/2$.  Hence
\cref{prop:gapped-deviation} holds with the explicit upper bound
\begin{equation}
    \epsGen(\Lambda,S,\delta)
    \leq
    \frac{4\Lambda(\Lambda+\Upsilon)}{\sqrt\mu}
    +3(\Lambda+\Upsilon)^2
      \sqrt{\frac{\log(8S/\delta)}{2\mu}}.
    \label{eq:geometric-generalization-bound}
\end{equation}
At fixed confidence, a raw-time gap of order $\log N$ is therefore sufficient
to keep the unit-block coupling penalty controlled under geometric mixing.
\end{corollary}

For a requested uniform-deviation tolerance $\epsilon_G$, abbreviate
\[
    a:=\frac4{\sqrt\mu},
    \qquad
    b:=3\sqrt{\frac{\log(4S/\delta_g)}{2\mu}},
\]
so that
$\epsGen(\Lambda,S,\delta)=(a+b)\Lambda^2+(a+2b)\Upsilon\Lambda+b\Upsilon^2$.

\begin{corollary}[Readout-capacity calibration]
\label{cor:ivanov-capacity-calibration}
If $\epsilon_G>b\Upsilon^2$, the unique positive radius satisfying
$\epsGen(\Lambda,S,\delta)=\epsilon_G$ is
\begin{equation}
    \Lambda
    =\frac{-(a+2b)\Upsilon
      +\sqrt{(a+2b)^2\Upsilon^2
        +4(a+b)(\epsilon_G-b\Upsilon^2)}}{2(a+b)}.
    \label{eq:ivanov-radius}
\end{equation}
If $\epsilon_G\leq b\Upsilon^2$, no positive Ivanov radius attains the
requested uniform tolerance with the declared $(N,g,S,\delta)$.
\end{corollary}

\paragraph{Why full-sample ERM remains controlled.}
The deviation is simultaneous over the entire finite pool and the entire
Ivanov ball.  Therefore the data-dependent minimizer in
\Eqref{eq:frozen-pool-erm} is already covered by the same event; no independent
validation fold is required for the theorem.  Searching more frozen designs
is charged through $\log S$, while fitting a richer readout is charged through
$\Lambda$.

\subsection{What qubit width can improve: the projection-information floor}
\label{sec:width-resolved-approximation}

A finite-data Johnson--Lindenstrauss statement cannot by itself imply lower
prediction risk.  The correct target-sensitive role of $n$ is the information
retained by $\projection_n$.  For a fixed realization of the coupled Gaussian
prefixes in \Eqref{eq:gaussian-projection}, let $\projection_n^{\sZ_-}$ act
coordinatewise on histories and define
\begin{equation}
    \projClass_n
    :=\left\{
      G\circ\projection_n^{\sZ_-}:
      G\in\gFFMP(\projection_n\gX)
    \right\},
    \qquad
    \projFloor_n
    :=\E\!\left[
      \left(
        F^\star(X_{-\infty:0})
        -\E\!\left[F^\star(X_{-\infty:0})\mid
          \sigma(\projection_nX_t:t\leq0)\right]
      \right)^2
    \right].
    \label{eq:projection-information-floor}
\end{equation}
This is the information-theoretic excess risk caused by observing only the
projected history.  Because $(\projection_n\gX)^{\sZ_-}$ is compact metrizable
and continuous functions are dense in $L^2$ of every Borel probability on it,
our FMP convention also gives the equivalent variational form
\begin{equation}
    \projFloor_n
    =\inf_{F\in\projClass_n}
      \E\!\left[(F(X_{-\infty:0})-F^\star(X_{-\infty:0}))^2\right].
    \label{eq:projection-floor-variational}
\end{equation}
Thus the floor isolates projection information loss before any restriction to
a particular reservoir realization.

\paragraph{Why the approximation proof may use $\reservoirUnitary_n=\identity$.}
The separation argument below uses the identity support point because it gives
an especially transparent universality certificate.  This should not be read
as a claim that mixing is redundant at finite resources.  The two questions
are logically different: \cref{sec:mixer-expressivity} proves that a generic
graph-local mixer exposes two- and three-stream interactions directly in raw
observable coordinates that the identity dynamics cannot expose at the same
observable locality, whereas the present section asks only whether the
support-closure family with an unrestricted universal readout has an intrinsic
approximation floor.  A universal readout can eventually compose local
coordinates even when the mixer is set to the identity; the mixer theorem
addresses how much interaction structure is already present before that
classical composition.

\begin{lemma}[Fading memory and projected-history separation]
\label{lem:feature-map-separation}
Fix $n\geq3$ and $\lambda_+\in(\lambda_0,1)$.  Every causal coordinate of an
admissible single-reservoir design is a continuous fading-memory functional of
the projected history.  Moreover, for every two distinct projected histories
$\vz^-\neq\vz'^{-}$ in $(\projection_n\gX)^{\sZ_-}$, there exists a
single-reservoir design in the support of the declared law whose causal
feature vector separates them.
\end{lemma}

\begin{theorem}[Width-resolved QuaRK approximation]
\label{thm:width-resolved-approximation}
Fix a realized preprocessor $\projection_n$ and
$\lambda_+\in(\lambda_0,1)$.  With
$\arch_R:=(n,R,\lambda_+)$,
\begin{equation}
    \lim_{\Lambda\to\infty}
      \inf_{R\geq1}
      \inf_{\design\in\operatorname{supp}(\designLaw_{\arch_R})}
      \epsApp(\design;\Lambda)
    =\projFloor_n.
    \label{eq:width-resolved-envelope}
\end{equation}
Equivalently, finite spatial multiplexing and finite-norm Mat\'ern readouts
have no asymptotic approximation floor beyond the information discarded by
$\projection_n$ itself.
\end{theorem}

\begin{corollary}[Monotone width and the injective regime]
\label{cor:width-information-monotonicity}
Under the canonical Gaussian coupling of \Eqref{eq:gaussian-projection},
\begin{equation}
    \projFloor_{n+1}\leq\projFloor_n,
    \qquad n\geq3.
    \label{eq:width-floor-monotonicity}
\end{equation}
If $n\geq d$, then $\projection_n$ has full column rank almost surely and
\begin{equation}
    \projFloor_n=0
    \qquad\text{almost surely}.
    \label{eq:width-floor-zero}
\end{equation}
No strict improvement is claimed at every increment of $n$: a target that
already factors through a narrower projection can have the same floor at
several widths.
\end{corollary}

\paragraph{Interpretation.}
This is the predictive role of width that the finite-data JL lemma alone cannot
supply.  Wider coupled projections never remove information available at a
narrower width, while $n\geq d$ eliminates the projection bottleneck for a
full-column-rank Gaussian map.  The theorem does not say that finite-sample
risk must decrease monotonically with $n$: the newly retained information may
be irrelevant to the target, and larger width still consumes more physical
qubits even though the RMS readout prevents coordinate count alone from
inflating the fidelity bounds.

\subsection{Good-design mass and finite-pool coverage}
\label{sec:pool-coverage}

The width-resolved theorem is an existence statement over the support.  To
connect it to random search, define the conditional probability mass of designs
whose causal approximation error lies within $\epsRep>0$ of the projection
floor,
\begin{equation}
    \goodMass_{\arch}(\epsRep;\Lambda\mid\projection_n)
    :=\designLaw_{\arch}\!\left(
       \epsApp(\design;\Lambda)
       \leq\projFloor_n+\epsRep
      \right).
    \label{eq:good-design-mass}
\end{equation}
The notation emphasizes that the design law is sampled after the shared
preprocessor has been frozen; the numerical value is target- and
projection-dependent.

\begin{proposition}[Support point implies positive good-design mass]
\label{prop:positive-good-mass}
Fix $\arch$, $\Lambda$, and $\epsRep>0$.  If some
$\design_0\in\operatorname{supp}(\designLaw_{\arch})$ satisfies
\begin{equation}
    \epsApp(\design_0;\Lambda)<\projFloor_n+\epsRep,
    \label{eq:strict-good-support-design}
\end{equation}
then
$\goodMass_{\arch}(\epsRep;\Lambda\mid\projection_n)>0$.
Consequently, by \cref{thm:width-resolved-approximation}, for every
$\epsRep>0$ there exist finite $R$ and finite $\Lambda$ for which the good-design
mass is positive.
\end{proposition}

\begin{proposition}[Finite-pool coverage]
\label{prop:finite-pool-coverage}
Let $\hitMass:=\goodMass_{\arch}(\epsRep;\Lambda\mid\projection_n)$.
Conditional on $\projection_n$, an i.i.d. pool of size $S$ satisfies
\begin{equation}
    \Prob\!\left[
      \min_{1\leq s\leq S}\epsApp(\design_s;\Lambda)
      >\projFloor_n+\epsRep
      \;\middle|\;\projection_n
    \right]
    =(1-\hitMass)^S.
    \label{eq:pool-miss-probability}
\end{equation}
\end{proposition}

\paragraph{Interpretation.}
The pool size now has a direct theoretical role: it suppresses the probability
of missing an $\epsRep$-good frozen representation exponentially fast.  The
paper proves positivity of $\hitMass$ when the architecture has a
strictly good support point, but it does not turn
\Eqref{eq:pool-miss-probability} into an a-priori numerical calibration of
$S$.  Such a calibration would require a useful quantitative lower bound on
$\hitMass$, which is target-, projection-, and architecture-dependent and is
not available under the present assumptions.  Campaign~III therefore measures
the empirical leverage of $S$ rather than treating a formal inversion of the
unknown hit mass as an operational prescription.

\subsection{ERM oracle inequality and end-to-end performance}
\label{sec:end-to-end-guarantee}

\begin{theorem}[Finite-pool ERM oracle inequality]
\label{thm:finite-pool-erm-oracle}
If $\delta_g>0$, then conditional on any realized $\projection_n$ and frozen
pool, with probability at least $1-\delta$ over the gapped sample and training
measurements,
\begin{equation}
    \oprisk_{\selectedDesign,w}(\selectedReadout)
    \leq
    \min_{1\leq s\leq S}
      \inf_{h\in\RKHS_{m_{\arch}}^\Lambda}
      \oprisk_{\design_s,w}(h)
    +2\epsGen(\Lambda,S,\delta).
    \label{eq:finite-pool-operational-oracle}
\end{equation}
\end{theorem}

The inequality is precisely the anti-overfitting role of the uniform bound:
the design and readout are selected on the full sample, but the price of doing
so is explicit and finite.

\begin{theorem}[Resource-to-performance certificate]
\label{thm:resource-performance}
Fix $\arch=(n,R,\lambda_+)$, $M,S,\Lambda$, and $\epsRep>0$, and set
$\hitMass:=\goodMass_{\arch}(\epsRep;\Lambda\mid\projection_n)$.
If $\delta_g>0$, then conditional on the realized shared projection, jointly
over the frozen pool, dependent sample, and training measurements, with
probability at least
\begin{equation}
    1-\delta-(1-\hitMass)^S,
    \label{eq:resource-performance-confidence}
\end{equation}
we have
\begin{equation}
\boxed{
    \oprisk_{\selectedDesign,w}(\selectedReadout)-\risk^\star
    \leq
    \projFloor_n+\epsRep
    +2\epsGen(\Lambda,S,\delta)
    +\epsShot(\arch,M)
    +\epsWin(\arch).
}
    \label{eq:resource-performance-bound}
\end{equation}
On the same event, the causal exact-feature counterpart obeys
\begin{equation}
    \risk_{\selectedDesign,\infty}(\selectedReadout)-\risk^\star
    \leq
    \projFloor_n+\epsRep
    +2\epsGen(\Lambda,S,\delta)
    +2\epsShot(\arch,M)
    +2\epsWin(\arch).
    \label{eq:resource-performance-causal}
\end{equation}
\end{theorem}

\paragraph{How to read the theorem.}
The terms have distinct roles.  Width $n$ controls the irreducible information
floor $\projFloor_n$; finite $R$ and $\Lambda$ determine whether the design law
places positive mass within $\epsRep$ of that floor; $S$ exponentially reduces
the probability of missing such a design but simultaneously appears through
$\log S$ in the statistical term; $M$ controls measurement fidelity; and
$(w,\lambda_+)$ control finite-window fidelity.  Unlike a conjunction of
unrelated certificates, every term in \Eqref{eq:resource-performance-bound}
participates in the risk of the predictor that the algorithm actually returns.

\begin{corollary}[Operational-budget calibration at fixed pool size]
\label{cor:resource-budget-calibration}
Fix tolerances $\epsRep,\epsilon_G,\epsilon_M,\epsilon_W>0$ and a declared
pool size $S$.  For fixed $\arch$, $\Lambda$, and realized $\projection_n$,
suppose $\hitMass>0$ and $\delta_g>0$.  If
\begin{equation}
\begin{aligned}
    \epsGen(\Lambda,S,\delta)&\leq\epsilon_G,\\
    M&\geq
    \left\lceil
      \frac{10\Lambda^2(\Lambda+\Upsilon)^2\kLip^2}{3\epsilon_M^2}
    \right\rceil,\\
    \lambda_+&\leq
    \min\!\left\{1,
      \left[\frac{\epsilon_W}
      {4\Lambda(\Lambda+\Upsilon)\kLip}\right]^{1/w}
    \right\},
\end{aligned}
    \label{eq:resource-budget-conditions}
\end{equation}
then, conditional on $\projection_n$, with probability at least
$1-\delta-(1-\hitMass)^S$,
\begin{equation}
    \oprisk_{\selectedDesign,w}(\selectedReadout)-\risk^\star
    \leq
    \projFloor_n+\epsRep+2\epsilon_G+\epsilon_M+\epsilon_W.
    \label{eq:calibrated-resource-performance}
\end{equation}
For fixed $(N,g,S,\delta)$, \cref{cor:ivanov-capacity-calibration} supplies the
corresponding admissible Ivanov radius for a requested $\epsilon_G$.  This
corollary calibrates the quantities whose bounds are explicit; it deliberately
does not prescribe $S$ from the unknown target-specific hit mass.
\end{corollary}

\subsection{Complementary finite-data preprocessing geometry}
\label{sec:finite-data-geometry}

The preceding theorem gives the predictive meaning of width.  A separate
finite-data question is how wide a Gaussian preprocessor must be to preserve
the geometry of the raw inputs that actually appear in the observed windows.
Let $\PointSet_N$ be their distinct values.  Since
$|\PointSet_N|\leq Nw$, the width can be chosen without inspecting their
realized geometry.

\begin{proposition}[Finite-data JL width]
\label{prop:finite-data-jl}
Fix $\epsilon_\Pi,\delta_\Pi\in(0,1)$.  If $Nw\geq2$ and
\begin{equation}
    n\geq n_{\mathrm{JL}}
    :=\frac{\log\!\left(2\binom{Nw}{2}/\delta_\Pi\right)}
      {\epsilon_\Pi^2/4-\epsilon_\Pi^3/6},
    \label{eq:joint-jl-width-rule}
\end{equation}
then, with probability at least $1-\delta_\Pi$ over $\projection_n$,
\begin{equation}
    (1-\epsilon_\Pi)\norm{u-v}_2^2
    \leq\norm{\projection_nu-\projection_nv}_2^2
    \leq(1+\epsilon_\Pi)\norm{u-v}_2^2
    \label{eq:joint-jl-geometry}
\end{equation}
for every distinct $u,v\in\PointSet_N$.  The statement is vacuous when
$Nw<2$.
\end{proposition}

The sufficient width is
$O(\epsilon_\Pi^{-2}\log(Nw/\delta_\Pi))$ and does not depend on the ambient
input dimension $d$.  This result is deliberately orthogonal to
\cref{thm:resource-performance}: it certifies the geometry of the observed raw
inputs, while $\projFloor_n$ quantifies what the projection can discard for the
population prediction problem.  Finite-data JL preservation is therefore not
used as a surrogate for approximation quality.
